\section*{Part V: Legal and Regulatory Compliance as an Operating Discipline}
\addcontentsline{toc}{section}{Part V: Legal and Regulatory Compliance as an Operating Discipline}

\subsection{31. Compliance is a calendar and an evidence system}

Maintain a single register of filings and renewals with requirement, entity, authority, legal basis, due date, accountable owner, external adviser, evidence of completion and escalation threshold. Include company secretarial filings, income tax, GST where applicable, TDS, payroll/labor, contracts, privacy, intellectual property, local premises permits and sector-specific education obligations. Actual requirements vary by the company and threshold.

A startup fails compliance when everyone assumes the accountant, lawyer or cofounder is handling it. Assign internal accountability even when work is outsourced.

\subsection{32. Companies Act and statutory records}

The Companies Act 2013 and applicable rules govern incorporation, director duties, accounts, audit, meetings, share issues and other corporate matters [G2]. Confirm whether the company qualifies for any small-company or startup-related procedural relief under current rules. Do not copy a compliance calendar from a much larger listed company without checking applicability; nor presume no annual filings are due because revenue is 0.

Keep statutory records consistent with bank financing, tax accounts and investor reports. Inconsistencies create due diligence and potential penalty risk.

\subsection{33. GST and income-tax distinctions}

Not every education service receives identical GST treatment. Classification can depend on provider status and the nature of services; a private training business should not assume exemptions applicable to qualifying educational institutions. Enterprise consulting, software licensing, overseas services and property transactions can have different tax consequences. Record invoice wording, place of supply, tax rate, input credits, withholding and collection rules with qualified advisers.

Pricing decisions should be based on net realizable revenue and after-tax cash, not a simplistic ``tuition is always tax-free'' claim.

\subsection{34. Employment and contractor governance}

Written agreements should cover role, compensation, benefits and statutory obligations, confidentiality, IP, acceptable use, conflicts, safeguarding, separation, and data access. A person described as a contractor may be treated differently under applicable law based on facts and control; labels alone do not resolve classification. Workers teaching minors or working in sensitive environments may trigger additional screening and safeguarding duties.

When international hiring becomes relevant, local employment, payroll, immigration and permanent-establishment rules must be separately reviewed.

\subsection{35. Data protection and security}

A technical institute handles learner identity, assessments, performance, payments, recordings, applicant information, and possibly enterprise client data. Establish legitimate purpose and permissions, minimization, role-based access, secure retention, vendor oversight, deletion/rights processes, incident response and protections for sensitive contexts. India's DPDP Act and 2025 Rules have staged implementation and subsequent official materials [G3]; verify provisions applicable on the action date rather than assuming every obligation began simultaneously.

Historical contacts of unclear consent status should not be used as a marketing list. A database creates responsibilities as well as apparent distribution power.

\subsection{36. Intellectual-property and employment conflict}

Founders still employed elsewhere should review employment agreements, outside-business policies, non-disclosure, invention assignment and conflict rules. Do not transfer employer code, client information, datasets or internal documents to the venture. Document independent development and use separate accounts and equipment where appropriate. No part of an employee's job description grants permission to exploit their prior employer's materials.

Obtain legal advice where overlapping research areas or invention ownership make the boundary uncertain.

\subsection{37. FDI and FEMA when investors are foreign}

A foreign investor may introduce permitted-sector, entry-route, pricing, reporting, bank KYC, beneficial ownership and downstream-investment obligations. RBI's foreign investment reporting framework includes the FIRMS platform and forms such as FC-GPR, FC-TRS and FLA in applicable situations [G4]. Do not rely on an old generic 30-day filing rule without confirming current regulations and transaction type.

An Indian founder should engage FEMA-qualified counsel and an authorized dealer bank before closing cross-border equity or convertible financing, not after money has arrived and reporting issues emerge.

\subsection{38. Foreign customer payments versus foreign donations}

An overseas learner's genuine fee for services and a foreign donor's contribution are not necessarily governed identically. India's Ministry of Home Affairs FCRA guidance distinguishes certain ordinary-course fees from foreign contribution and specifies registration/prior-permission requirements for eligible recipients of foreign contributions [G5]. The classification depends on facts and current law. A sponsor cannot relabel a donation as a commercial invoice to avoid regulation.

Maintain separate agreements, bank accounts and restricted-fund records as required. Foreign philanthropy deserves advance specialist advice.

\subsection{39. CSR, grants, philanthropy and restrictions}

Indian CSR funding can require eligible implementing agencies and prescribed filings. Grants and donations may be restricted to scholarships, research, equipment or infrastructure. The recipient cannot treat a restricted grant as ordinary shareholder profit. Track eligibility, donor representations, use of proceeds, reporting, conflicts and refunds. A founder-controlled commercial enterprise and a future nonprofit entity must not mingle restricted and unrestricted cash.

An endowment requires separate investment and spending rules, authorized fiduciaries and donor-respecting governance.

\subsection{40. Education and degree recognition}

UGC explains that recognized degrees derive from qualifying statutory universities or other lawful degree-granting routes, subject to standards and applicable councils [G6]. A company providing a serious non-degree technical course cannot simply call itself a university or promise degree equivalence. Use precise labels for certificates, professional learning, partner-issued awards and credits, and obtain legal review of collaboration claims.

Formal academic expansion requires separate sponsor, curriculum, campus, faculty and statutory approval work. The commercial company's registration is not university authorization.

\subsection{41. Maharashtra private-university pathway}

The Maharashtra Private Universities (Establishment and Regulation) Act, 2023, updated official text and the 2024 Rules are relevant to a possible state private university [G7]. They describe a proposal and scrutiny process, academic and infrastructure requirements and a statutory architecture. The state may amend its law or interpret requirements for specific sponsors and locations. A founder should commission a compliance gap analysis and DPR before purchasing land for a hoped-for university.

The board must see the legal pathway, actual sponsor eligibility, asset holding rules, financing and academic leadership - not a generic PowerPoint saying ``approval is possible.''
